\section{Related Work}\label{sec:related}

\paragraph{Multi-turn attention degradation.}
Recent work has established that LLMs systematically degrade across conversation turns. \citet{levy2025lost} showed that model performance drops on tasks placed in later turns, even when the tasks are identical to earlier ones. \citet{liu2026intent} traced a related failure to intent mismatch, where models progressively shift from the user's original intent as dialogue extends. MT-Bench-101 \citep{bai2024mtbench101} provides fine-grained multi-turn evaluation across 1388 turn-level instances, revealing that capability varies substantially across task categories. \citet{dongre2026attention} offered a mechanistic account: they defined the Goal Attention Ratio (GAR) and showed that attention to system-prompt goal tokens decays by over 70\% across turns in 12 models, while residual stream probes retain goal information with AUC\,=\,0.98. This dissociation between attention collapse and residual persistence motivates our gating intervention, and our analysis reveals why such interventions fail to transfer from training to inference.

\paragraph{Residual gating and adaptive depth.}
Controlling computation flow through residual connections has a rich history. ReZero \citep{bachlechner2021rezero} showed that initializing residual multipliers at zero enables stable training of very deep networks, establishing the principle that residual scaling can shape optimization dynamics. Subsequent work extended this idea to variable-depth inference: LayerDrop \citep{fan2020layerdrop} applies structured dropout to entire layers during training, Mixture-of-Depths \citep{raposo2024mixture} routes tokens past layers via learned per-token decisions, FlexiDepth \citep{luo2025flexidepth} adds adaptive layer-skipping to pretrained LLMs, and CALM \citep{schuster2022calm} implements early exit based on confidence scores. These methods share the insight that not all layers contribute equally to every input, but they target computational efficiency or training stability. ThreadGate applies residual gating to a different objective: recovering attention patterns in multi-turn dialogue, a setting where adaptive depth methods have not been evaluated.

\paragraph{Attention manipulation.}
A parallel line of work modifies attention distributions directly rather than through the residual pathway. StreamingLLM \citep{xiao2024streaming} identified the attention sink phenomenon, where initial tokens absorb disproportionate attention regardless of content, and proposed retaining anchor tokens to stabilize streaming inference. \citet{ko2024found} showed that middle-positioned information receives systematically less attention and proposed positional calibration to correct this bias. \citet{qiu2025gated} introduced gating within the attention computation itself, replacing softmax with a learned non-linear gate that produces sparse, sink-free attention patterns. These approaches intervene inside the attention mechanism; ThreadGate instead operates on the residual stream. Our results suggest that neither pathway alone overcomes the distributional gap between teacher-forced training and autoregressive inference.

\paragraph{Multi-task gradient conflict.}
When multiple loss terms compete for shared parameters, their gradients can interfere. The multi-task optimization literature addresses this through gradient manipulation: GradNorm \citep{chen2018gradnorm} normalizes gradient magnitudes across tasks, PCGrad \citep{yu2020pcgrad} projects away conflicting components, CAGrad \citep{liu2021cagrad} seeks descent directions that reduce worst-case task loss, and Nash-MTL \citep{navon2022nash} frames task balancing as Nash bargaining. All four methods treat gradient conflict as a root cause to be resolved through better optimization. We repurpose PCGrad as a causal analysis tool: by eliminating directional conflict between NTP and GAR gradients and observing that GAR convergence does not improve, we establish gradient conflict as a symptom of the training dynamics rather than the cause of the failure.

\paragraph{Exposure bias.}
Teacher forcing creates a distributional mismatch between training and inference that compounds over sequence length. Scheduled Sampling \citep{bengio2015scheduled} introduced curriculum-based mixing of ground-truth and model-generated tokens. MIXER \citep{ranzato2016sequence} combined token-level and sequence-level objectives, and Professor Forcing \citep{lamb2016professor} used adversarial regularization to align hidden state distributions under both regimes. \citet{he2019exposure} showed that autoregressive models exhibit self-recovery from exposure bias at the token prediction level. SEARNN \citep{leblond2018searnn} proposed locally recurrent search-based training to avoid the mismatch entirely. This literature focuses on generation quality degradation in primary objectives. Our work identifies a distinct failure mode: auxiliary attention losses that converge under teacher forcing fail completely under autoregressive generation, even when NTP remains robust, because token-level self-recovery does not extend to distributional properties of attention.